\chapter{Baseline Inference and the Discovery of a Sensitivity Failure: From Scene-Level to Person-Relative Grounding (Day 3)}
\label{chap:day03}

\section{Motivation}

With a fixed, balanced, rule-mapped sample of 163 images in hand from Day 2, the pilot's next
step was to produce a baseline judgment for every sample: does the model say each image complies
with the safety rule it was assigned, or not? Every later evaluation day --- masking-based
Descriptive Accuracy, attention-based explanation extraction, Robustness, Bounded Completeness ---
depends on this baseline answer being a meaningful judgment in the first place, one that actually
discriminates between compliant and violating images rather than defaulting to the same answer
regardless of input. Day 3's task was therefore not only to run inference over the sample, but to
verify, quantitatively, that the resulting proxy was fit for that purpose before any explanation
method was built on top of it. As this chapter documents, that verification surfaced a serious
problem in the most natural implementation of the proxy, and resolving it became the day's central
piece of work.

\section{Implementation}

Florence-2-base-ft has no native free-form visual-question-answering task token: its task
vocabulary consists of grounding and captioning tokens --- \texttt{<CAPTION>}, \texttt{<OD>},
\texttt{<OPEN\_VOCABULARY\_DETECTION>}, \texttt{<CAPTION\_TO\_PHRASE\_GROUNDING>} --- rather than a
head designed to answer arbitrary yes/no questions \citep{xiao2024florence2}. Each of the four
safety rules therefore had to be reframed as a grounding query rather than posed as a direct
question. \texttt{inference.py}'s \texttt{answer\_rule()} dispatches on the rule being tested:
rule~4 (struck-by risk, worker-to-excavator proximity) is routed to \texttt{\_answer\_rule\_4},
which performs a proximity check between two independently grounded objects, \texttt{"worker"} and
\texttt{"excavator"}; rules~1--3 (hard hat, harness, guardrail) are routed to
\texttt{\_answer\_presence\_rule}, which --- in its first implementation --- simply asks whether the
relevant safety object can be grounded anywhere in the image at all.

\texttt{scripts/03\_run\_baseline\_inference.py} ran this logic over all 163 samples from Day 2's
\texttt{pilot\_samples.csv}, writing one row per sample to \texttt{results/baseline\_predictions.csv}
and rendering 20 overlay images, with detected boxes drawn on the original frame, to
\texttt{results/figures/baseline/}.

\section{Architecture Before and After}

Figure~\ref{fig:diagram-b} contrasts the two versions of the rules~1--3 proxy built and run on Day
3. The left-hand side, v1, is the scene-level existence check described above: a single
\texttt{<OPEN\_VOCABULARY\_DETECTION>} call for the safety object, and an answer that depends only on
whether that call returns any box at all. The right-hand side, v2, is the person-relative redesign
developed later the same day in response to the failure documented in
Section~\ref{sec:day03-problems}: two independent detection calls, one for \texttt{"worker"} and one
for the safety object, followed by a per-worker coverage check rather than a single scene-wide
existence check. Both versions share the same rule~4 proximity logic, which is not shown again here
because it was not changed.

\begin{figure}[h]
\centering
\begin{tikzpicture}[node distance=1.1cm]
\node[block] (img1) {Image};
\node[block, below=of img1] (ovd1) {\texttt{<OVD>}\\(safety object)};
\node[block, below=of ovd1] (exist1) {Box exists\\anywhere?};
\node[block, below=of exist1, fill=rulegray] (ans1) {Answer:\\compliant / violation};
\draw[arrow] (img1) -- (ovd1);
\draw[arrow] (ovd1) -- (exist1);
\draw[arrow] (exist1) -- (ans1);
\node[diagcaption, above=0.1cm of img1] {\textbf{v1: scene-level} (3.0\% sensitivity)};

\node[block, right=3.2cm of img1] (img2) {Image};
\node[block, below left=1.1cm and 0.3cm of img2] (worker2) {\texttt{<OVD>}\\("worker") $\to$\\worker\_boxes};
\node[block, below right=1.1cm and 0.3cm of img2] (obj2) {\texttt{<OVD>}\\(safety object) $\to$\\object\_boxes};
\node[block, below=2.3cm of img2] (cover2) {all\_boxes\_covered\\(per-worker check)};
\node[block, below=of cover2, fill=rulegray] (ans2) {Answer:\\compliant / violation};
\draw[arrow] (img2) -- (worker2);
\draw[arrow] (img2) -- (obj2);
\draw[arrow] (worker2) -- (cover2);
\draw[arrow] (obj2) -- (cover2);
\draw[arrow] (cover2) -- (ans2);
\node[diagcaption, above=0.1cm of img2] {\textbf{v2: person-relative} (27.0\% sensitivity, 9x)};
\end{tikzpicture}
\caption{Safety-rule proxy logic before (left) and after (right) the Day 3 redesign.}
\label{fig:diagram-b}
\end{figure}

\section{Worked Examples}
\label{sec:day03-worked-examples}

Before quantifying the proxy's accuracy, the pipeline's basic mechanics were validated by visual
inspection of the rendered overlays. Figure~\ref{fig:day03-rule1} shows image 7: a rule~1 (hard
hat) query, where \texttt{<OPEN\_VOCABULARY\_DETECTION>} returns a plausible, tightly-localized box
landing correctly on a worker's head. The scene itself contains five workers, a detail that becomes
important in Section~\ref{sec:day03-limitations}'s discussion of the residual single-worker-box
limitation. Figure~\ref{fig:day03-rule4} shows image 120: rule~4's proximity check, where both
\texttt{"worker"} and \texttt{"excavator"} are independently grounded, and the resulting boxes stand
in a spatially sound relationship to one another. Together, these two examples confirm that the
underlying detections are real and well-localized, not noise --- a fact that turns out to matter
directly in Section~\ref{sec:day03-problems}, where the v1 proxy's failure is traced to how the
detections are \emph{used}, not to any defect in the detections themselves.

\begin{figure}[h]
\centering
\includegraphics[width=0.85\linewidth]{baseline/0000007_rule_1.png}
\caption{Image 7, rule~1 (hard hat): \texttt{<OPEN\_VOCABULARY\_DETECTION>} correctly localizes a
hard hat on a worker's head. Five workers are visible in the original scene.}
\label{fig:day03-rule1}
\end{figure}

\begin{figure}[h]
\centering
\includegraphics[width=0.85\linewidth]{baseline/0000120_rule_4.png}
\caption{Image 120, rule~4 (struck-by risk): the proximity check grounds both \texttt{"worker"} and
\texttt{"excavator"} with a sound spatial relationship between the two boxes.}
\label{fig:day03-rule4}
\end{figure}

Pipeline mechanics held up at the full sample size, not only on these two hand-picked examples: all
163 target \texttt{image\_id}s from Day 2's sample were found in the dataset stream, and
\texttt{assigned\_rule\_id} round-tripped correctly for every row. Mean inference time was 174ms per
call and mean confidence (mean token probability) was 0.58, both consistent with Day 1's
single-sample numbers and indicating that nothing changed about the model's basic operating
characteristics when moving from one image to a 163-image batch.

\section{Results}

\subsection{v1: scene-level proxy}

With the pipeline confirmed mechanically sound, v1's predictions were compared against ground-truth
\texttt{primary\_class} to measure sensitivity (the fraction of true violations the proxy correctly
flags as violations) and specificity (the fraction of true compliant images the proxy correctly
flags as compliant). The presence-based proxy covering rules~1--3 and the proximity-based proxy
covering rule~4 were scored separately, since they are structurally different checks.

\begin{center}
\begin{tabular}{lrr}
\toprule
\textbf{Proxy (v1)} & \textbf{Sensitivity} & \textbf{Specificity} \\
\midrule
Presence-based (rules 1--3) & 3.0\% (3/100) & 97.4\% (37/38) \\
Proximity-based (rule 4) & 61.5\% (8/13) & 41.7\% (5/12) \\
\bottomrule
\end{tabular}
\end{center}

The contrast between these two rows is stark and is the central finding of Day 3. The
presence-based proxy achieves high specificity but almost no sensitivity at all: of 100 true
violations across rules~1--3, it correctly flags only 3. A proxy with 3.0\% sensitivity is, for
practical purposes, a near-constant "compliant" classifier --- it agrees with the dataset on
compliant images almost every time, and disagrees on violations almost every time, which is the
signature of a classifier that is not discriminating on the relevant signal at all. The
proximity-based rule~4 proxy, by contrast, already shows real, balanced discriminative signal at
this small sample size (62\% sensitivity, 42\% specificity on $n=13$ and $n=12$ respectively) and
required no redesign.

\subsection{v2: person-relative proxy}

After the redesign described in Section~\ref{sec:day03-problems}, \texttt{scripts/03\_run\_baseline\_inference.py}
was re-run on all 163 samples with the updated \texttt{\_answer\_presence\_rule} logic. The results:

\begin{center}
\begin{tabular}{lrr}
\toprule
\textbf{Proxy (v2)} & \textbf{Sensitivity} & \textbf{Specificity} \\
\midrule
Presence-based (rules 1--3), overall & \textbf{27.0\% (27/100)} & 73.7\% (28/38) \\
\quad rule\_1 (hard hat) & 16.0\% (8/50) & --- \\
\quad rule\_2 (harness) & 30.8\% (4/13) & --- \\
\quad rule\_3 (guardrail) & 40.5\% (15/37) & --- \\
Proximity-based (rule 4), unchanged & 61.5\% (8/13) & 41.7\% (5/12) \\
\bottomrule
\end{tabular}
\end{center}

Overall presence-based sensitivity rises from 3.0\% to \textbf{27.0\%} --- a \textbf{9x
improvement} --- at the cost of specificity falling from 97.4\% to 73.7\%. This is an expected, and
in this case a desirable, precision/recall tradeoff: the v1 proxy's high specificity was an
artifact of a classifier that almost never says "violation," not evidence of genuine discriminative
skill, and the v2 proxy's lower specificity is the necessary consequence of a proxy that has
started to actually discriminate instead of defaulting to "compliant." Per-rule sensitivity is
uneven: rule~3 (guardrail) benefits most, rising to 40.5\%, plausibly because its looser proximity
threshold (described below) is closer in spirit to the relational rule~4 check that already worked
well; rule~1 (hard hat) improves least, to 16.0\%, though even this is more than five times its v1
value.

\section{Problems Encountered and Resolutions}
\label{sec:day03-problems}

The 3.0\% sensitivity figure for the v1 presence-based proxy was the one substantial problem
encountered on Day 3, and diagnosing its root cause was the day's central piece of analytical work.

\textbf{Root cause.} This was not a code bug: the underlying detections are real, well-localized
boxes, a fact already confirmed visually in Section~\ref{sec:day03-worked-examples}'s worked
examples. The failure is a mismatch between what the dataset's ground-truth label means and what
the v1 query actually asks. The dataset's \texttt{rule\_1\_violation} field (and the equivalent
fields for rules~2 and~3) marks that \emph{a specific worker} in the image lacks the required PPE
or protection. \texttt{<OPEN\_VOCABULARY\_DETECTION>}, as used in v1, answers a \emph{scene-level
existence} question instead --- "does a hard hat appear anywhere in this image?" Most
ConstructionSite images contain multiple workers. If even one worker in the frame is wearing a hard
hat, the existence query returns a box, and the v1 proxy answers "compliant," even when a different
worker in the very same frame is the one the dataset actually flagged as the violation.

\textbf{Concrete traced case.} Image 79 (ground truth \texttt{fall\_hazard}, rule~3, missing edge
protection) makes the mechanism concrete. The "guardrail" query matched the wooden formwork
bordering the excavation pit --- a structure that is guardrail-shaped but is not a guardrail ---
and so the proxy answered "compliant," even though the dataset's actual violation is a missing
barrier elsewhere in the same scene. The detection itself was not wrong in any technical sense: a
guardrail-shaped object genuinely is present in the image. The query was simply answering a
different, easier question than the one the ground-truth label was actually recording.

\textbf{Why rule~4 was unaffected.} The rule~4 proximity proxy does not share this failure mode,
because it is relational rather than existential: it asks about the spatial relationship between
two specific grounded objects (worker and excavator), not whether a safety object exists anywhere
in the scene. This is consistent with rule~4 already showing real, balanced discriminative signal
(62\% sensitivity, 42\% specificity) before any redesign was attempted.

\textbf{Resolution, stated as a decision.} Rather than silently patch around individual failure
cases such as image 79, the rules~1--3 proxy was redesigned, as an explicit and documented decision
rather than an incidental bug fix, to a \textbf{person-relative} check. The redesigned
\texttt{\_answer\_presence\_rule} detects all \texttt{"worker"} boxes and all relevant safety-object
boxes independently, then flags \textbf{violation} if any one detected worker has no nearby or
overlapping safety-object box, via a new \texttt{regions.all\_boxes\_covered} helper. This
mirrors the dataset's own semantics directly: one non-compliant worker is enough for the dataset to
mark the whole image violated, so the proxy now requires \emph{every} detected worker to be covered
before it will answer "compliant," rather than requiring only that the safety object exist
somewhere in the frame. Worn-PPE rules (rule~1, hard hat; rule~2, harness) use a tight proximity
threshold, \texttt{PRESENCE\_PROXIMITY\_FRACTION = 0.08} (8\% of the image's maximum dimension),
reflecting that the safety object must be on the worker's body to count as worn. Rule~3's
guardrail is structural and can protect a worker from further away, so it uses a looser threshold,
0.20 (20\%). Images with zero detected workers --- 12 of the 163 samples --- fall back to the
original v1 scene-level existence check, since the person-relative logic has no worker box to
anchor to in those cases.

\section{Standard vs. Non-Standard Choices}

Two distinct design decisions are at work in this chapter, and they sit at opposite ends of the
standard/non-standard spectrum. \textbf{Standard}: treating a yes/no safety-compliance question as
a closed-vocabulary grounding or detection task is a documented workaround pattern for
vision-language models that, like Florence-2, lack a native free-form VQA head --- reframing the
question in terms of the model's actual task vocabulary, rather than forcing an unsupported task
onto it, follows established practice for this class of model.

\textbf{Non-standard}, and the chapter's central methodological contribution: the person-relative
redesign of \texttt{\_answer\_presence\_rule} itself. This is not a textbook XAI or VQA technique
drawn from the literature; it is a domain-specific adaptation that mirrors the ConstructionSite
dataset's own per-worker labeling semantics --- one non-compliant worker marks the whole image
violated --- and encodes that semantics directly into the proxy's logic via independent worker and
object detection plus a per-worker coverage check. It was arrived at only after the v1 proxy's
near-total sensitivity failure was diagnosed and traced to a specific, concrete case (image 79), not
designed in advance. We report it here as a decision, with its quantified before/after effect, rather
than folding it silently into the baseline as if it had always been the design.

\section{Limitations}
\label{sec:day03-limitations}

The person-relative redesign is bounded by a residual limitation in Florence-2's detection behavior
that the redesign itself cannot work around. \texttt{<OPEN\_VOCABULARY\_DETECTION>} queried with the
singular phrase \texttt{"worker"} returns only \emph{one} box, even in scenes that visibly contain
several workers. This was confirmed visually on two of the sample's own images: image 7 (five
workers visible, ground truth \texttt{ppe\_violation}, shown in Figure~\ref{fig:day03-rule1}) and
image 79 (multiple people visible, ground truth \texttt{fall\_hazard}, the same image traced in
Section~\ref{sec:day03-problems}). When the single worker box that Florence-2 happens to return is
the compliant one, the per-worker coverage check has nothing to flag, even though the dataset's
ground-truth violation concerns a different, undetected worker in the same frame. This caps how far
the person-relative redesign can push sensitivity using the current detection strategy: the v2 proxy
can only be as discriminating as the worker boxes it is given, and a single-box-per-scene detection
will always miss multi-worker violations where the detected worker is not the violating one.

A different detection strategy --- for example, \texttt{<DENSE\_REGION\_CAPTION>}, which is designed
for exhaustive instance enumeration rather than single best-match grounding --- could plausibly
close this residual gap by returning all worker instances rather than one. This was deliberately
\emph{not} pursued in this pass: the 9x sensitivity gain already achieved by the person-relative
redesign is a substantial, honestly-earned improvement in its own right, and chasing the residual
gap immediately risked conflating two separate questions (is the person-relative logic correct, and
is the underlying detection exhaustive) within a single change. The gap is documented here instead
as an explicit, bounded limitation for future work.

\section{Implications for the Paper}

The headline number for this chapter, and the one a paper reviewer will most want explained
clearly, is the 3.0\%-to-27.0\% (9x) sensitivity improvement in the rules~1--3 proxy, achieved by
recognizing that the dataset's ground truth is a per-worker judgment and redesigning the proxy to
match that judgment rather than continuing to ask a scene-level existence question. This redesign
is not incidental to the pilot's later chapters: it is a precondition for them. Descriptive
Accuracy, Robustness, and Bounded Completeness, evaluated in later chapters, all assume that the
Day~3 baseline answer is a meaningful judgment worth testing an explanation method against. Under
the \emph{old} v1 proxy, for rules~1--3, that assumption was nearly false --- a near-constant
"compliant" classifier gives a masking- or perturbation-based explanation method nothing to flip,
because there is almost no correct "violation" judgment in the first place for an explanation to be
faithful to. This is precisely why the redesign was carried out on Day 3, before any of those later
evaluation days ran, rather than discovered and patched retroactively after building explanation
methods on top of a proxy that could not discriminate.

For the eventual paper, we recommend a dedicated "Threats to Validity" paragraph addressing the
residual detection gap documented in Section~\ref{sec:day03-limitations}: that Florence-2's
open-vocabulary detection for the singular phrase \texttt{"worker"} returns at most one instance per
scene, and that this bounds the person-relative proxy's achievable sensitivity independently of how
well its per-worker coverage logic is designed. Reporting this limitation alongside the 9x
improvement allows the improvement to be read accurately --- as a genuine, substantial gain that is
nonetheless capped by an instance-enumeration constraint in the underlying vision-language model,
rather than as a fully solved problem. With a baseline proxy now genuinely discriminating between
compliant and violating images, the pilot's next step is to extract and evaluate the explanations
that accompany these judgments.
